\section*{الفصل \ref{ch:g10:biodiversity-scales} --- التنوع الحيوي على كل سلّم}

\begin{solution}{exo:g10:biodiversity-scales:1}
\omterm{def:g10:biodiversity-scales:biodiversity}{الأنظمة البيئية} (بركة، غابة زان)؛ \omterm{def:g10:biodiversity-scales:species}{والأنواع} (ضفادع البركة وسمادلها
وقصبها)؛ \omterm{def:g10:universal-dna:gene}{والمورّثات} (\omterm{def:g10:universal-dna:gene}{أليلات} \omterm{def:g10:universal-dna:gene}{مورّثة} ABO في جمهرة بشرية).
\end{solution}

\begin{solution}{exo:g10:biodiversity-scales:2}
جماعة من الكائنات المتشابهة تتناسل فيما بينها وتعطي نسلًا خصبًا.
والأسد والنمر يعطيان هجائن عقيمة، فهما نوعان.
\end{solution}

\begin{solution}{exo:g10:biodiversity-scales:3}
$\num{1000000}/6400 \approx 156$: أي نحو 150 ضعفًا.
\end{solution}

\begin{solution}{exo:g10:biodiversity-scales:4}
نسبة نسخ \omterm{def:g10:universal-dna:gene}{مورّثة} ما في جمهرة التي هي \omterm{def:g10:universal-dna:gene}{أليل} معين. وB $= 1 - 0.66 - 0.27 = 0.07$.
\end{solution}

\begin{solution}{exo:g10:biodiversity-scales:5}
خمسة؛ وأكبرها انقراض نهاية البرمي، قبل 252 مليون سنة، وقد أزال نحو
95\% من \omterm{def:g10:biodiversity-scales:species}{الأنواع} البحرية.
\end{solution}

\begin{solution}{exo:g10:biodiversity-scales:6}
الثاني: فالغنى واحد، لكن الوفرة النسبية أكثر توازنًا بكثير. والغنى
وحده لا يفرق بينهما؛ وإنما تفرق الوفرات النسبية.
\end{solution}

\begin{solution}{exo:g10:biodiversity-scales:7}
الحقل نفسه؛ والزميل عاين عشرة أضعاف. \omterm{def:g10:biodiversity-scales:species}{والأنواع} الموجودة تزداد مع الجهد
حتى يستوي المنحني، فالأعداد المأخوذة بجهد مختلف غير قابلة للمقارنة ---
فالعدد 8 عند التلميذ معاينة ناقصة، لا حقل أفقر.
\end{solution}

\begin{solution}{exo:g10:biodiversity-scales:8}
الناجون القلائل لم يحملوا إلا جزءًا من \omterm{def:g10:universal-dna:gene}{أليلات} \omterm{def:g10:biodiversity-scales:species}{النوع}؛ وذريتهم، مهما
كثرت، تتقاسم تلك المجموعة المختزلة. \omterm{def:g10:biodiversity-scales:biodiversity}{فالتنوع الوراثي} يفقد في عنق
الزجاجة ولا ترده الأعداد: فليس لدى \omterm{def:g10:biodiversity-scales:species}{النوع} احتياط في وجه مرض جديد أو
مناخ متغير (وحال الفهد الصياد مثال).
\end{solution}

\begin{solution}{exo:g10:biodiversity-scales:9}
نحو 560 عائلة قبله و300 بعده: أي $260/560 \approx 45\%$ من العائلات
مفقودة. والعائلة لا تختفي إلا باختفاء أنواعها كلها، فقد نجت عائلات
كثيرة \omterm{def:g10:biodiversity-scales:species}{بنوع} أو نوعين وفقدت الباقي؛ ولذلك كانت نسبة \omterm{def:g10:biodiversity-scales:species}{الأنواع} المفقودة
(نحو 95\%) أكبر بكثير من نسبة العائلات.
\end{solution}

\begin{solution}{exo:g10:biodiversity-scales:10}
كل جمهرة تحمل \omterm{def:g10:universal-dna:gene}{الأليلات} الثلاثة، ويختلف أفرادها فيما بينهم فيما يحملون؛
أما الفروق بين الجمهرات فليست إلا إزاحات في التواتر. ومن ثم فمعظم
\omterm{def:g10:biodiversity-scales:biodiversity}{التنوع الوراثي} \omterm{def:g10:biodiversity-scales:species}{للنوع} موجود داخل كل جمهرة: فجاران يختلفان عادة في عدد
من المواضع كالذي يختلف فيه شخصان متباعدان.
\end{solution}

\begin{solution}{exo:g10:biodiversity-scales:11}
\omterm{def:g10:biodiversity-scales:biodiversity}{النظام البيئي}: البركة نفسها بمجتمعها. \omterm{def:g10:biodiversity-scales:species}{والنوع}: السمندل أو الرعّاش الذي
لم يكن يعيش محليًا إلا فيها. \omterm{def:g10:universal-dna:gene}{والمورّثات}: \omterm{def:g10:universal-dna:gene}{الأليلات} التي حملتها جمهرة
ضفادع البركة، وقد تكون مختلفة عن \omterm{def:g10:universal-dna:gene}{أليلات} برك أخرى.
\end{solution}

\begin{solution}{exo:g10:biodiversity-scales:12}
بعمر يبلغ نحو $\num{2e6}$ سنة، تعطي $\num{2e6}$ \omterm{def:g10:biodiversity-scales:species}{نوع} نحو انقراض واحد
في السنة. وألف ضعف ذلك نحو 1000 \omterm{def:g10:biodiversity-scales:species}{نوع} في السنة، أي $\num{100000}$ في
القرن --- أي 5\% من \omterm{def:g10:biodiversity-scales:species}{الأنواع} المعروفة في مئة سنة.
\end{solution}

\begin{solution}{exo:g10:biodiversity-scales:13}
أخلى الانقراض أنماط عيش --- عاشبات كبيرة، وضواري كبيرة، وطائرات،
وسابحات --- كانت الديناصورات تشغلها. فتكيفت سلالات الثدييات الناجية،
وقد تحررت من أولئك المنافسين ولها تنوعها الوراثي تسحب منه، مع الأدوار
الشاغرة وانقسمت إلى \omterm{def:g10:biodiversity-scales:species}{أنواع} جديدة: فالخسارة فتحت الفرص التي ملأها
التنوع الجديد.
\end{solution}

\begin{solution}{exo:g10:biodiversity-scales:14}
فقد القمح المزروع معظم \omterm{def:g10:universal-dna:gene}{أليلات} أسلافه. أما الأقارب البرية فتحفظها:
\omterm{def:g10:universal-dna:gene}{أليلات} لمقاومة مرض، أو جفاف، أو ملوحة، لا يحملها أي صنف حالي.
وإدخالها في القمح بالتهجين يرد الخيار؛ وحالما ينقرض قريب بري تضيع
أليلاته إلى الأبد.
\end{solution}

\begin{solution}{exo:g10:biodiversity-scales:15}
الانقراض طبيعي، لكن بنحو \omterm{def:g10:biodiversity-scales:species}{نوع} واحد في السنة؛ أما النسبة الحاضرة فأسرع
مئات إلى ألف مرة، وهي تضاهي انقراضًا جماعيًا، والتعافي من تلك
الانقراضات استغرق ملايين السنين --- أي أطول مما وجد نوعنا. وفي الوقت
نفسه، \omterm{def:g10:biodiversity-scales:biodiversity}{فالأنظمة البيئية} التي تفقد تقدم الغذاء والتأبير والتربة وتنقية
الماء والاحتياطات الوراثية لمحاصيلنا. فالسابقة تبين أن الخسارة حقيقية
وأن التعافي غير متاح على أي سلّم زمني بشري.
\end{solution}

\begin{solution}{pb:g10:biodiversity-scales:1}
\textbf{1.} 31 و9 \omterm{def:g10:biodiversity-scales:species}{أنواع}: والنسبة $31/9 \approx 3.4$.

\textbf{2.} المزرعة A: لا \omterm{def:g10:biodiversity-scales:species}{نوع} فوق 22\% مقابل \omterm{def:g10:biodiversity-scales:species}{نوع} عند 84\% في المزرعة
B. والتوازن يقول كيف تتوزع الأفراد على \omterm{def:g10:biodiversity-scales:species}{الأنواع} الموجودة؛ فالمجتمع
الغني غير المتوازن يسوده \omterm{def:g10:biodiversity-scales:species}{نوع} واحد ويكون الباقي نادرًا.

\textbf{3.} الزيادات تتقلص (7 و4 و2): فالمنحني يستوي، والعينة قريبة من
التمام؛ وبضعة مربعات أخرى لن تضيف أكثر من \omterm{def:g10:biodiversity-scales:species}{نوع} أو نوعين.

\textbf{4.} بلغ منحني المزرعة B العدد 9 عند المربع الخامس عشر ولم يضف
شيئًا بعده: فمنحني المعاينة الكاملة يستوي إلى هضبة، وقد بلغتها
المزرعة B وكادت تبلغها المزرعة A.

\textbf{5.} كثير من النباتات لا يرى ولا يعرف إلا وهو مزهر، والحشرات
والطيور تتغير مع الموسم؛ فشهر آخر كان ليغير الأعداد لأسباب لا صلة لها
بالسياج.

\textbf{6.} الحشرات $46/14 \approx 3.3$؛ والطيور $17/5 = 3.4$.
فالمجموعات الثلاث تعطي عاملًا يقارب 3: أي القصة نفسها.

\textbf{7.} السياج يوفر النباتات (أوراقًا وأزهارًا وخشبًا ميتًا
ومأوى) التي تتغذى عليها الحشرات وتتكاثر فيها؛ \omterm{def:g10:biodiversity-scales:species}{وأنواع} نباتية أقل تعني
\omterm{def:g10:biodiversity-scales:species}{أنواع} حشرات أقل؛ والحشرات وثمار السياج ومواضع تعشيشه تطعم الطيور
وتؤويها، فيهبط تنوع الطيور بدوره.

\textbf{8.} لا: فالمحصول لا يقيس إلا الزرع. والسياج قد يقي الحقل من
الريح، ويمسك التربة، ويؤوي الملقحات وضواري آفات الزرع، وآثاره تظهر في
السنوات السيئة أو عبر عقود، لا في رقم موسم واحد.

\textbf{9.} طائر يعشش في السياج كالشحرور أو دخلة السياج (أو حلزون
السياج، أو بقة الزعرور الدرعية): فهو يحتاج إلى الشجيرات للتعشيش أو
الغذاء أو المأوى، وهي ما لم تعد المزرعة B تقدمه.

\textbf{10.} \omterm{def:g10:biodiversity-scales:biodiversity}{الأنظمة البيئية}: فالسياج نفسه، وهو موئل، قد أزيل، وذهبت
معه \omterm{def:g10:biodiversity-scales:species}{الأنواع} \omterm{def:g10:universal-dna:gene}{والمورّثات}؛ أما الخسارات على السلمين الآخرين فنتائج لذلك.

\textbf{11.} المزرعة A: نسخ y $2 \times 100 + 200 = 400$، ونسخ b
$200 + 2 \times 100 = 400$، من أصل 800: فتواتر y هو 0.5. والمزرعة B:
y $8 + 32 = 40$، وb $32 + 128 = 160$، من أصل 200: فتواتر y هو 0.2.

\textbf{12.} المزرعة A: الأليلان متساويا التواتر و50\% من الحلزونات
تحمل كليهما؛ والمزرعة B: 32\% تحمل كليهما وb هو السائد. فالمزرعة A
أكثر تنوعًا.

\textbf{13.} على الضفة العارية تؤكل الحلزونات الصفراء أكثر؛ فجيلًا بعد
جيل انتقلت نسخ y أقل، وهبط تواتر y من نحو 0.5 إلى 0.2.

\textbf{14.} تنعكس الأفضلية: فتبقى الحلزونات الصفراء أفضل في السياج
المظلل ويعود تواتر y إلى الارتفاع عبر الأجيال --- شرط أن يكون \omterm{def:g10:universal-dna:gene}{الأليل}
لا يزال موجودًا.

\textbf{15.} حلزونات المزرعة B: ففي جمهرة من بضع مئات، تستطيع سنة
سيئة أو بضع وفيات بالمصادفة أن تزيل كل نسخة من \omterm{def:g10:universal-dna:gene}{الأليل} الأندر؛ أما بين
عشرات الألوف \omterm{def:g10:universal-dna:gene}{فالأليل} تحمله آلاف الحلزونات ولا تقدر المصادفة على
محوه.

\textbf{16.} المزرعة A: $31 + 46 + 17 = 94$؛ والمزرعة B:
$9 + 14 + 5 = 28$؛ والنسبة $94/28 \approx 3.4$.

\textbf{17.} ثلاثة بالمئة من الأرض تحوي موئل نحو ثلثي \omterm{def:g10:biodiversity-scales:species}{أنواع} المزرعة.

\textbf{18.} النباتات $24/31 \approx 77\%$؛ والحشرات
$35/46 \approx 76\%$؛ والطيور $12/17 \approx 71\%$.

\textbf{19.} \omterm{def:g10:biodiversity-scales:biodiversity}{التنوع الوراثي}: \omterm{def:g10:biodiversity-scales:species}{فالأنواع} تعود من جمهرات مجاورة، أما
\omterm{def:g10:universal-dna:gene}{الأليلات} التي فقدتها الجمهرات المحلية (حلزونات المزرعة B مثلًا) فلا
تعود إلا إذا صادف أن حملها وافدون؛ \omterm{def:g10:universal-dna:gene}{والأليل} المفقود لا تعيد إعادة
الغرس خلقه.

\textbf{20.} ضاعف السياج النباتي غنى \omterm{def:g10:biodiversity-scales:species}{أنواع} المزرعة نحو 3.4 مرة في كل
مجموعة عدت؛ والسياج المعاد غرسه يرد \omterm{def:g10:biodiversity-scales:biodiversity}{النظام البيئي}، ومع الوقت معظم
\omterm{def:g10:biodiversity-scales:species}{الأنواع}، لكنه لا يرد وحده \omterm{def:g10:universal-dna:gene}{الأليلات} التي فقدتها الجمهرات المحلية.
\end{solution}
